\section{Overview}

In Viotti et al. \cite{viotti2016consistency}, authors provided a
structured overview
of different consistency semantics appeared in distributed systems. Building
upon it's semantic definitions, we propose to use predicate logic derived from
those definitions to model consistency semantics, and reason about the
compositional properties of these semantics with SMT solvers (while Z3
  \cite{demoura2008z3} is used in our implementation, it's easily
  extensible with other
solver backends like CVC5 \cite{barbosa2022cvc5}).

To capture behaviors in distributed system executions, we model state
transitions on individual nodes using the notion of
operations. Similar to
Viotti et al. \cite{viotti2016consistency}, we define an operation as a tuple
containing: \texttt{proc} (process or node), \texttt{type} (read or
write), \texttt{obj} (the object
being acted upon), \texttt{ival} (input value), \texttt{oval} (output
value), \texttt{stime}
(start time), and \texttt{rtime} (return time). In our SMT encoding,
this tuple is
modeled as a sort, serving as the atomic building block for capturing system
behaviors and state transition constraints.  Additional fields can be added
to this sort by invoking the sort constructor
with extra parameters, adding the capability to model more complex
operations if needed.
The instantiations of operation tuple represents atomic actions
in the system.

Axioms and relations are predicate logic constraints defined over operations.
Partial and total orderings of operations are assigned through relation
functions defined in history and abstract
execution \cite{burckhardt2014principles}
\cite{viotti2016consistency}. In the implementation, they are defined
as uninterpreted functions with constraints \cite{bryant2002modeling}.
History represent the set of all operations executed in the system,
while abstract execution instantiates/refines
the history by specifying which operations are visible to, or are arbitrated
over each other. By embedding these relations in first order logic, we can
leverage SMT solvers to check the satisfiability of a given configuration,
thereby determining whether a particular consistency model holds or if certain
axioms of that model are violated. Axioms defined in history are
deterministic orderings based on time and session
(returns-before, same-session, session-order, etc.) \needcite{}. Whereas
abstract execution relations are non-deterministic orderings
used to capture more complex temporal constraints and conflict
resolution (visibility, arbitration, happens-before, etc.) \needcite{}.

Consistency semantics or models are properties of a system or
composition of systems
defined as logical formulas over axioms/relations (e.g., Monotonic Reads, PRAM,
Causal Consistency, etc.). Each model enforces safety properties i.e.
specific ordering
and visibility constraints on operations within history and abstract
executions. % TODO: add more

On top of our defined consistency models, we can perform individual consistency
semantic specification and verification, as well as pairwise checks between
two models to determine if one model's execution trace is statisfiable under
another model's constraints. Though this is described as a pairwise check,
each model can be the combination of multiple consistency models.
This allows us to reason about the interaction
between different consistency models, and whether a safety property of one model
can be used to satisfy or violate the safety property of another
model. Compositional
checks can be performed by construction a multigraph of systems
(nodes) connected by
edges (data flow and constraints between systems) and our tool can
output result of
all satisfiable assignments for each system in the graph.

Compatibility (pairwise) check (\texttt{compatible(lhs, rhs)})
determines if one set of constraints (lhs) implies another set of
constraints (rhs) under a given background context (others and global
Relation.Constraints()). This function uses the construct helper
function to create a Z3 solver instance. construct asserts the
negation of the implication: Not(Implies(lhs, rhs)). This is
logically equivalent to lhs And Not(rhs). It also adds the background
constraints (others and Relation.Constraints()). The compatible
function then calls s.check() on this solver. If s.check() returns
z3.unsat (unsatisfiable), it means there is no model (no
counterexample) where lhs is true and rhs is false. Therefore,
whenever lhs holds, rhs must also hold. The function returns True
(compatible). If s.check() returns z3.sat (satisfiable), it means a
counterexample exists, the implication does not hold universally, and
the function returns False (not compatible).
